\documentclass{article}
\usepackage[utf8]{inputenc}
\usepackage{amsmath}
\usepackage{geometry}
\geometry{margin=2cm}
\begin{document}

\section*{Análise Detalhada da Questão 174 — ENEM 2024 — Matemática}

\subsection*{Enunciado}
Em um jogo virtual para celular, um personagem pode percorrer trajetórias retilíneas voando ou se deslocando ao longo de paredes. Considere que o personagem descreve a trajetória \( ABCDEF \), em que os pontos \( A, D \) e \( E \) estão em um plano paralelo ao que contém os pontos \( B \) e \( C \), sendo esses dois planos ortogonais ao plano da base que contém o ponto \( F \), conforme a figura.

A projeção ortogonal, sobre o plano da base, da trajetória \( ABCDEF \) descrita pelo personagem é:

\subsection*{Solução Técnica}

A projeção correta é dada pela alternativa \textbf{C}. Para obtê-la, projetamos cada ponto da trajetória verticalmente para o plano da base, anulando sua coordenada de altura, e conectamos os pontos na ordem dada.

\bigskip

\textbf{Passos principais:}
\begin{enumerate}
  \item Identificar a posição dos planos e pontos.
  \item Eliminar a coordenada vertical \( z \) para obter a projeção ortogonal.
  \item Conectar os pontos \( A, B, C, D, E, F \) na sequência correta.
  \item Comparar com as alternativas e concluir que a alternativa C está correta.
\end{enumerate}

\bigskip

\textbf{Fórmula de projeção:}
\[
(x,y,z) \to (x,y,0)
\]

\subsection*{Explicação Didática}

Para compreender a projeção ortogonal, imagine a "sombra" que a trajetória faz ao ser iluminada por uma luz vertical na direção \( z \). Projeta-se cada ponto diretamente para baixo no plano base.

Devemos:
\begin{itemize}
  \item Conhecer bem os planos envolvidos e as posições dos pontos.
  \item Anular a altura \( z \) durante a projeção.
  \item Observar a ordem correta dos pontos ao ligar a trajetória.
\end{itemize}

\subsection*{Erros Comuns}
\begin{itemize}
  \item Trocar a ordem dos pontos ao projetar.
  \item Manter a coordenada de altura na projeção.
  \item Confundir os pontos \( B \) e \( C \).
  \item Projetar numa direção diferente da vertical.
  \item Achar que a projeção deve preservar inclinação.
\end{itemize}

\subsection*{Dicas para Ensino}
\begin{itemize}
  \item Usar software ou maquetes para visualização tridimensional.
  \item Relacionar projeção ortogonal a "sombras".
  \item Enfatizar análise espacial e ordenação dos pontos.
  \item Explicar que a projeção reduz dimensões.
\end{itemize}

\subsection*{Análise dos Distratores}

\textbf{Alternativa A:} Confusão na ordem da ligação ou na direção da projeção, aparentando similaridade parcial, mas incorreta na conexão.

\textbf{Alternativa B:} Confusão entre posições dos pontos \( B \) e \( C \), com ordem errada e conexão incorreta.

\textbf{Alternativa D:} Não elimina corretamente a altura de \( D \) e \( E \), distorcendo a trajetória no plano.

\textbf{Alternativa E:} Desconhecimento dos planos e projeção correta, formando trajetória que não corresponde à real projeção ortogonal.

\subsection*{Perfil da Questão}

Foca na habilidade de interpretação espacial e projeção ortogonal, requerendo visualização 3D e análise geométrica sem cálculos numéricos complexos. O nível é médio, adequado para o ENEM.

\bigskip

Em resumo, a correta interpretação dos planos e a aplicação da projeção ortogonal evidenciam a alternativa \textbf{C} como a resposta correta para a questão.

\end{document}